\item Identifica problemas de la vida cotidiana y plantea soluciones.\\[-1.2em]
\item Conoce y caracteriza el pensamiento científico para plantearse y resolver problemas en la escuela y su cotidianidad.\\[-1.2em]
\item Valora la influencia del conocimiento científico y tecnológico en la sociedad actual.\\[-1.2em]
\item Identifica las unidades de medición que se ocupan en su entorno escolar, familiar y en su comunidad.\\[-1.2em]
\item Identifica cuáles son, cómo se definen y cuál es la simbología de las unidades básicas y derivadas del Sistema Internacional de Unidades.\\[-1.2em]
\item Realiza conversiones con los múltiplos y submúltiplos al referirse a una magnitud.\\[-1.2em]
\item Conoce los instrumentos de medición, materiales, sus propiedades y características.\\[-1.2em]
\item Relaciona e interpreta las teorías sobre estructura de la materia, a partir de los modelos atómicos y de partículas y los fenómenos que les dieron origen.\\[-1.2em]
\item Explora algunos avances recientes en la comprensión de la constitución de la materia y reconoce el proceso histórico de construcción de nuevas teorías.\\[-1.2em]
\item Experimenta e interpreta los modelos atómicos y de partículas al proponer hipótesis que expliquen los tres estados de la materia, sus propiedades físicas como la temperatura de fusión, ebullición, densidad, entre otros.\\[-1.2em]
\item Interpreta la temperatura y el equilibrio térmico con base en el modelo de partículas.\\[-1.2em]